%% related_work.tex

\section{Related Work}
\label{sec:related-work}

We position this work along three axes: empirical software engineering methods for codebase analysis, architectural analyses of LLM agent frameworks, and trust assumption research in self-healing systems.

\subsection{Empirical Software Engineering and Codebase Analysis}

Empirical software engineering has a long tradition of source-level code analysis to identify systematic patterns and deficiencies across codebases. Runeson and H{\"o}st~\cite{runeson2009} established the methodological foundations for case study research in software engineering, including the multi-case design and validity threat classification adopted in this paper. Kitchenham et al.~\cite{kitchenham2002} provided guidelines for empirical research that combine qualitative and quantitative evidence. Wohlin et al.~\cite{wohlin2012} systematized the classification of validity threats (internal, external, construct, conclusion) that we apply in \S\ref{sec:threats-to-validity}.

Mining software repositories (MSR) is a complementary methodology that uses automated tools to extract patterns from version control histories and issue trackers~\cite{bird2015}. Our approach differs from typical MSR studies in that we perform \emph{manual} source-level analysis with a written qualitative codebook adapted from established coding practice~\cite{glaser1967,strauss1998}; we do not claim a complete Grounded Theory study. This is appropriate for our research question, which concerns a structural pattern (the trust-escalation trampoline) that requires semantic understanding of the recovery path's trust assumptions, not just pattern matching on code tokens. Storey~\cite{storey2008} reviews similar qualitative approaches to software documentation and code analysis.

\subsection{Architectural Analysis of LLM Agent Frameworks}

The LLM agent framework landscape has been surveyed from multiple perspectives. Reflexion~\cite{reflexion} introduced verbal reinforcement learning as a self-healing pattern; AutoGen~\cite{autogen} proposed multi-agent conversation as an architectural paradigm; SWE-agent~\cite{swe-agent} demonstrated agent-computer interfaces for automated software engineering. MetaGPT~\cite{meta-agent} formalized multi-agent collaboration through role assignments; LangGraph~\cite{langgraph} provided a graph-based orchestration framework; Aider~\cite{aider-code} and OpenHands~\cite{openhands-code} focused on coding agent workflows. These works describe the \emph{functional} architecture of their frameworks but do not systematically analyze the \emph{trust assumptions} embedded in their recovery paths. Our work complements these architectural descriptions by providing an empirical analysis of a cross-cutting concern (provenance handling) that spans all six frameworks.

\subsection{Trust Assumptions in Self-Healing Systems}

The security literature has studied prompt injection on the forward execution path~\cite{greshake2023indirect,liu2024formalizing,willison2023prompt}, and recent work has proposed defenses including instruction hierarchy enforcement~\cite{instruction-hierarchy}, spotlighting~\cite{spotlighting}, and capability-based control flow separation~\cite{debenedetti2025camel}. Zhan et al.~\cite{zhan2025adaptive} demonstrated that adaptive attacks can bypass forward-path defenses, achieving $>$50\% ASR against eight proposed defenses. We engage directly with Zhan et al.'s threat model to position \rci{} relative to indirect prompt injection. Both share the common root cause of untrusted external content entering a high-trust context; however, \rci{} constitutes a distinct attack surface that existing forward-path defenses do not cover, on three grounds.

First, the \emph{injection channel} differs. Indirect injection---including Zhan et al.'s adaptive variants---injects content via the tool's \emph{return value}, read directly by the LLM on the forward path. \rci{} injects content via the tool's \emph{failure}, captured in the error trace and re-injected by the recovery pipeline. The two channels are governed by different code paths: forward-path defenses that inspect tool return values do not inspect error traces (\S\ref{sec:defense-aware}).

Second, the \emph{authority transition} differs. Indirect injection relies on the LLM obeying content that arrives in a tool message of standard trust level. \rci{} relies on the recovery pipeline placing untrusted content into an action-authorizing retry or reflection context (the trust-escalation trampoline, RQ1 in \S\ref{sec:method-taxonomy}). This transition---performed by the framework, not the attacker---is the software path that our taxonomy documents across four of six frameworks; it need not literally use a system-role message in every case.

Third, the \emph{defense-bypass mechanism} differs. Zhan et al. show that
adaptive attackers can craft payloads that evade forward-path defenses by
content manipulation. \rci{} instead enters through a different framework
path: the adversarial text is embedded in failure content after the forward
input has already been accepted. The author-constructed five-class screen
tests this channel separately and is reported as supplementary evidence in
\S\ref{sec:defense-aware}; it is not pooled with the native framework
results. Recovery-path provenance tracking is therefore a distinct design
requirement, even when a system already deploys forward-path controls.

These works primarily target the forward path. We did not find an empirical software-engineering study that isolates the \emph{recovery path} as a separate trust boundary across agent frameworks.

Classical operating systems enforce trust boundaries via privilege levels and provenance tags~\cite{saltzer1975protection,levine2003system,bovet2005understanding}. Our work shows that a classical trust-boundary vulnerability pattern reappears systematically in self-healing LLM agent frameworks: the recovery pipeline is a privileged component (it re-injects content into the LLM context under a high-trust frame) that acts on failure content of unknown provenance. The OWASP Top 10 for LLM Applications~\cite{owasp-llm-top10} enumerates forward-path injection risks but does not classify the recovery channel as a distinct trust boundary. Our empirical analysis suggests it should be.

This paper provides cross-framework empirical evidence of the recovery-path trust-boundary pattern, together with a controlled evaluation of content-level provenance tagging on a native recovery path.
